\documentclass[notes=show]{beamer}
\usepackage{mathpazo}
\usepackage{hyperref}
\usepackage{multimedia}
\usepackage{tcolorbox}
\usepackage{tikz}
\usetikzlibrary{shadows}
\usetheme{metropolis}
\usecolortheme{default}
\setbeamertemplate{itemize item}{%
    \begin{tikzpicture}
        \shade[ball color=black!100!yellow, preaction={fill=black,
opacity=.00,transform canvas={xshift=1mm,yshift=-1mm, yscale=0.5}}] (0,0) circle (0.6ex);
    \end{tikzpicture}
}
\setbeamertemplate{itemize subitem}{%
    \begin{tikzpicture}
        \shade[ball color=black!100!white, preaction={fill=black,
        opacity=.00,transform canvas={xshift=1mm,yshift=-1mm, yscale=0.5}}] (0,0) circle (0.6ex);
    \end{tikzpicture}
}
\setbeamercolor{frametitle}{bg=white}
\setbeamercolor{frametitle}{fg=black}
\setbeamercolor{background canvas}{bg=white}
\setbeamercolor{block body}{bg=mDarkTeal!30}
\setbeamercolor{block title}{bg=mDarkTeal,fg=black!2}
\setbeamertemplate{navigation symbols}{}
\setbeamertemplate{footline}[page number]
\newenvironment{stepenumerate}{\begin{enumerate}[<+->]}{\end{enumerate}}
\newenvironment{stepitemize}{\begin{itemize}[<+->]}{\end{itemize} }
\newenvironment{stepenumeratewithalert}{\begin{enumerate}[<+-| alert@+>]}{\end{enumerate}}
\newenvironment{stepitemizewithalert}{\begin{itemize}[<+-| alert@+>]}{\end{itemize} }

\begin{document}

\title[The Skill Content]{The Skill Content of Recent Technological Change: An Empirical Exploration}
\subtitle{}
\date{David H. Autor, Frank Levy \& Richard J. Murnane\bigskip \\
\textit{Quarterly Journal of Economics}, November 2003}
\author{}
\maketitle

\section{Introduction}

\begin{frame}{What explains skill-biased technical change?}
\begin{itemize}
\item There is a positive correlation between the adoption of computers and relative demand for college workers.\medskip
\item This has been interpreted as evidence of Skill-Biased Technical Change (SBTC). \medskip
\item However, this interpretation labels the correlation without explaining its cause.\medskip
\item Why do computers shift relative labor demand?
\end{itemize}
\end{frame}

\begin{frame}{Preview of findings}
\begin{itemize}
\item Computerization will reduce demand for routine labor tasks and increase demand for non-routine labor tasks, and workers reallocate from routine tot non-routine tasks (I.A). \medskip
\item These effects are stronger in industries that are initially more intensive in routine labor tasks (I.B). \medskip
\item The paper tests these predictions using task measures (II-V). \medskip
\item It shows that task changes due to computerization explain a substantial fraction of the increase in relative demand for college workers (VI).
\end{itemize}
\end{frame}

\begin{frame}{Table of Contents}
\begin{itemize}
\item[I.] The Task Model \smallskip
\item[II.] Empirical Implementation \smallskip
\item[III.] Trends in Job Task Input, 1960-1998 \smallskip
\item[IV.] Computerization and Tasks: Industry-Level Relationships \smallskip
\item[V.] Task Change Within Education Groups and Occupations \smallskip
\item[VI.] Quantifying the Magnitude of Task Structure Changes \smallskip
\item[VII.] Conclusions
\end{itemize}
\end{frame}

\section{I. The Task Model}

\subsection{I.A The demand for routine and non-routine tasks}

\begin{frame}{I. The task model}
\begin{enumerate}
\item[\textcolor{red}{I.A}] \textcolor{red}{The demand for routine and non-routine tasks} \bigskip
\item[\textcolor{gray}{I.B}] \textcolor{gray}{Industry-level implications}
\end{enumerate}
\end{frame}

\begin{frame}{Routine and non-routine tasks}
\begin{itemize}
\item Technological progress is seen as routine labor tasks than can be automated.\medskip
\item Routine labor tasks follow explicit rules that can be written in computer code.\medskip
\item Non-routine labor tasks cannot be specified in computer code.\medskip
\item Labor task can also be analytic and interactive or manual.
\end{itemize}
\end{frame}

\newpage
\begin{center}
\begin{figure}
\includegraphics[height=8.5cm,width=\textwidth]{table1.png}
\end{figure} 
\end{center}
\newpage

\begin{frame}{Intuitive assumptions}
Consider the following 3 assumptions: \medskip
\begin{enumerate}
\item Routine and non-routine tasks are imperfect substitutes. \medskip
\item Computers substitute for routine labor tasks.\medskip
\item Computers increase the productivity of non-routine labor tasks.
\end{enumerate} \medskip
The paper develops a formal model based on these assumptions.
\end{frame}

\begin{frame}{Production function}
\begin{itemize}
\item Assume an aggregate Cobb-Douglas production function:
\[
Q=[L_{R}+C]^{1- \beta}L_{N}^{ \beta} \text{ with } \beta \in (0,1)  \tag{1}  \label{eq1}
\]
with $L_{R}$ and $L_{N}$ routine and non-routine labor and $C$ computer capital, in efficiency units. \bigskip
\item $C$ is supplied perfectly elastically at price $ \rho $. \bigskip
\item $ \rho $ is falling exogenously with time due to technological progress. This is the causal force in the model.
\end{itemize}
\end{frame}

\begin{frame}{Assumptions (formally)}
This production function implies the following 3 assumptions:
\begin{enumerate} \medskip
\item The elasticity of substitution between $L_{R}$ or $C$ and $L_{N}$ is $1$, i.e. routine and non-routine tasks are imperfect substitutes. \medskip
\item $C$ substitutes for routine labor tasks because the elasticity of substitution between $C$ and $L_{R}$ is $\infty $. \medskip
\item The marginal productivity of non-routine labor tasks rises with computerization because $C$ and $L_{N}$ are q-complements.
\end{enumerate}
\end{frame}

\begin{frame}{Labor supply: Roy framework}
\begin{itemize}
\item Each worker $i$ inelastically supplies one unit of labor. \medskip
\item Workers have heterogeneous productivity endowments in both routine and non-routine labor tasks:
\[
\forall i: E_{i}=[r_{i},n_{i}] \text{ with } 1 \geqslant r_{i}, n_{i} > 0
\]
\item Worker $i$ chooses $ 0 \leqslant \lambda_{i} \leqslant 1 $ to maximize income:
\[
L_{i}=[\lambda_{i}r_{i},[1 - \lambda_{i}]n_{i}]
\]
\item Workers will choose a task based on comparative advantage as in Roy [51] and some will re-allocate following a decline in $ \rho $.
\end{itemize}
\end{frame}

\begin{frame}{Equilibrium}
\begin{itemize}
\item The wage per efficiency unit of routine labor tasks must be equal to the price of computer capital:
\[
w_{R} = \rho  \tag{2}  \label{eq2}
\]
\item Define $ \eta_{i} \equiv n_{i}/r_{i} \in (0, \infty)$ as $i$'s relative efficiency at doing non-routine versus routine labor tasks. \medskip
\item Define the relative efficiency of marginal workers as:
\[
\eta^{*} = w_{R}/w_{N}  \tag{3}  \label{eq3}
\]
\item Worker $i$ supplies routine labor tasks ($ \lambda_{i}=1 $) if $ \eta_{i} < \eta^{*} $ and supplies non-routine labor tasks otherwise ($ \lambda_{i}=0 $).
\end{itemize}
\end{frame}

\begin{frame}{Equilibrium: Sorting}
\begin{itemize}
\item Labor supply in efficiency units to routine and non-routine labor tasks is given by:
\[
g( \eta ) = \sum_{i} r_{i} I( \eta_{i}  < \eta ) \text{ and } h( \eta ) = \sum_{i} n_{i} I (\eta_{i} \geqslant \eta) 
\]
with $I$ an index function. \medskip
\item Further assume that $ \eta_{i} $ has nonzero support at all $ \eta_{i} \in (0, \infty )$, so that $g( \eta )$ is continuously upward sloping and $h( \eta )$ is continuously downward sloping in $\eta$. \medskip
\item This ensures that $L_{R}^{*}$ and $L_{N}^{*}$ are not zero.
\end{itemize}
\end{frame}

\begin{frame}{Equilibrium: Productive efficiency}
\begin{itemize}
\item Productive efficiency requires:
\[
w_{R} = \frac{ \partial Q}{ \partial L_{R}}=[1- \beta] \theta^{-\beta} \text{ and } w_{N} = \frac{ \partial Q}{\partial L_{N}}= \beta \theta^{1-\beta} \tag{4}  \label{eq4}
\]
where $ \theta $ is the ratio of routine to non-routine tasks:
\[
\theta \equiv [C+g(\eta^{*})]/h( \eta^{*}) \tag{5}  \label{eq5}
\]
\item Equations (2) to (5) provide equilibrium conditions for the model's variables $(w_{R},w_{N}, \theta, \eta)$. \medskip
\item How does a decline in the price of computers affect task inputs, relative wages, and labor supply?
\end{itemize}
\end{frame}

\begin{frame}{Technological progress, relative wages, and labor supply}
\begin{itemize}
\item From the expression for $w_{R} = \rho$ in equation (4) we get that:
\[
\frac{\partial ln(\theta)}{\partial ln(\rho)} = - \frac{1}{\beta} < 0
\]
\item Combining this with equations (2)-(4) we get that:
\begin{tcolorbox}
\[
\frac{\partial ln (w_{N}/w_{R})}{\partial ln(\rho)} = - \frac{\partial ln(\eta^{*})}{\partial ln(\rho)} = \frac{\partial ln(\theta)}{\partial ln(\rho)} = - \frac{1}{\beta} < 0 \tag{7}  \label{eq7}
\]
\end{tcolorbox}
\item The relative wage of workers doing non-routine tasks increases and marginal workers will reallocate from routine to non-routine labor tasks.
\end{itemize}
\end{frame}

\begin{frame}{Summary so far}
\begin{itemize}
\item An exogenous decline in the price of computers increases the demand for computers. \medskip
\item The demand for routine labor tasks declines because computers and routine labor tasks are perfect substitutes. \medskip
\item The demand for non-routine labor tasks increases because routine and non-routine tasks are q-complements. \medskip
\item Marginal workers reallocate their labor supply from routine tot non-routine labor tasks.
\end{itemize}
\end{frame}

\subsection{I.B Industry-level implications}

\begin{frame}{I. The task model}
\begin{enumerate}
\item[\textcolor{gray}{I.A}] \textcolor{gray}{The demand for routine and non-routine tasks} \bigskip
\item[\textcolor{red}{I.B}] \textcolor{red}{Industry-level implications}
\end{enumerate}
\end{frame}

\begin{frame}{Industry-level implications}
\begin{itemize}
\item An aggregate production function and decline in the price of computers only gives a single time series to test the model. \medskip
\item Additional leverage comes from assuming different sectors $j$ with different task technologies, i.e. with different $ \beta_{j}$. \medskip
\item The more routine-task intensive an industry initially is ($ \beta_{j}$ small), the more will be invested in computers and the larger will be the decline in routine labor tasks and the rise in non-routine labor tasks.
\end{itemize}
\end{frame}

\begin{frame}{Industry-level implications}
\begin{itemize}
\item Assume industry $j$'s production function is:
\[
q_{j} = r_{j}^{1- \beta_{j}} n_{j}^{ \beta_{j}} \text{ with } \beta_{j} \in (0,1) \tag{8}  \label{eq8}
\]
\item Assume Dixit-Stiglitz monopolistic competition with utility of the representative consumer given by:
\[
U(q_{1},...,q_{j},...) = [ \sum_{j} q_{j}^{1- \nu}]^{1/[1- \nu]} \text{ with } 0 < \nu <1  \tag{9}  \label{eq9}
\]
\item Consumer demand for each good $j$ is given by:
\[
q_{j} = [p_{j}/P^{1 - \nu}]^{-1/ \nu} Q
\]
with $P$ the aggregate price index and $p_{j} = MC/[1 - \nu ]$. 
\end{itemize}
\end{frame}

\begin{frame}{Industry factor demands}
\begin{itemize}
\item Profit maximization by firms gives demands for routine and non-routine tasks in each industry: \medskip
\item Demand for routine tasks is given by:
\[
r_{j}= \rho^{-1/ \nu} [[1- \beta_{j}][1- \nu]]^{1/ \nu} \left[ \frac{w_{N}}{\rho} \frac{[1- \beta_{j}]}{\beta_{j}} \right]^{[\beta_{j}[ \nu - 1]] / \nu} \tag{11a}  \label{eq11}
\]
\item Demand for non-routine labor tasks is given by:
\[
n_{j}=w_{N}^{-1/ \nu} [ \beta_{j}[1- \nu]]^{1/ \nu} \left[ \frac{w_{N}}{\rho} \frac{[1- \beta_{j}]}{\beta_{j}} \right]^{[[1- \beta_{j}][1 - \nu]] / \nu} \tag{11b}  \label{eq11}
\]
\end{itemize}
\end{frame}

\begin{frame}{Industry-level predictions}
\begin{enumerate}
\item The proportionate increase in routine tasks is larger in routine-task intensive industries ($ \beta_{j} $ small) (P1):
\begin{tcolorbox}
\[
\frac{ \partial ln(r_{j})}{ \partial \rho} = \frac{ \beta_{j}[1 - \nu] - 1}{\nu \rho} <0 \text{ and } \frac{\partial ^2 ln(r_{j})}{\partial \rho \partial \beta_{j}} = \frac{1- \nu}{\nu \rho} >0
\]
\end{tcolorbox}
\item The proportionate increase in non-routine tasks is larger in routine-task intensive industries (P2):
\begin{tcolorbox}
\[
\frac{\partial ln(n_{j})}{\partial \rho} = \frac{[\beta_{j}-1][1 - \nu]}{\nu \rho} <0 \text{ and } \frac{\partial ^2 ln(n_{j})}{\partial \rho \partial \beta_{j}} = \frac{1 - \nu}{\nu \rho} >0
\]
\end{tcolorbox}
\item Similar predictions can be made for $j$ an occupation. (P3)
\end{enumerate}
\end{frame}

\section{II. Empirical Implementation}

\subsection{II.A Selecting measures of routine and non-routine tasks}

\begin{frame}{II. Empirical implementation}
\begin{enumerate}
\item[\textcolor{red}{II.A}] \textcolor{red}{Selecting measures of routine and non-routine tasks} \bigskip
\item[\textcolor{gray}{II.B}]  \textcolor{gray}{A predictive test} \bigskip
\end{enumerate}
\end{frame}

\begin{frame}{Main data}
\begin{itemize}
\item \underline{Task measures:} Fourth [1977] Edition and Revised Fourth [1991] Edition of Dictionary of Occupational Titles (DOT) examine 12K occupations along 44 dimensions.\medskip
\item \underline{Worker-level data:} Census 1\% extracts for 1960, 1970, 1980 and 1990 and CPS (MORG) files for 1980, 1990, 1998. Variables are FTE hours worked per year, gender, education, 485 Census Occupation Codes (COC) and 140 industries. \medskip
\item Changes in labor tasks have an extensive margin (i.e. changes in occupational employment holding task content within occupations at the 1977 DOT level) and an intensive margin (i.e. changes in labor tasks within occupations comparing the 1977 and 1991 DOT).
\end{itemize}
\end{frame}

\newpage
\begin{center}
\begin{figure}
\includegraphics[height=8.5cm,width=\textwidth]{appendixtable1.png}
\end{figure} 
\end{center}
\newpage

\newpage
\begin{center}
\begin{figure}
\includegraphics[height=8.5cm,width=\textwidth]{DOT77.png}
\end{figure} 
\end{center}
\newpage

\begin{frame}{DOT task measures: 1960 centiles}
\begin{itemize}
\item The DOT measures don't have a natural scale. \medskip
\item Each DOT measure is transformed into percentile values corresponding to their occupational rank weighted by occupational employment using the 1960 Census. \medskip
\item Task measures are relative to the 1960 distribution of tasks across workers, measured in centiles. \medskip
\item Not normalizing the raw DOT scores gives results that are qualitatively identical.
\end{itemize}
\end{frame}

\subsection{II.B A predictive test}

\begin{frame}{II. Empirical implementation}
\begin{enumerate}
\item[\textcolor{gray}{II.A}] \textcolor{gray}{Selecting measures of routine and non-routine tasks} \bigskip
\item[\textcolor{red}{II.B}]  \textcolor{red}{A predictive test} \bigskip
\end{enumerate}
\end{frame}

\begin{frame}{A predictive test using the merged DOT and 1960 Census}
\begin{itemize}
\item Define routine tasks across occupations $m$ in industry $j$ as:
\[
r_{j} \equiv \sum_{m} [STS_{jm} + FINGDEX_{jm}]
\]
\item Define the sum of all 5 task inputs $k$ in industry $j$ as:
\[
r_{j} + n_{j} \equiv \sum_{m} \sum_{k} T_{jmk}
\]
\item The percentage Routine Task Share for industry $j$ then is:
\[
RTS_{j,1960}=100 \times r_{j,1960}/[r_{j,1960} + n_{j,1960}]
\]
\end{itemize}
\end{frame}

\begin{frame}{A predictive test using the merged DOT and 1960 Census}
\begin{itemize}
\item To measure computer adoption, use CPS to calculate an industry's percentile rank of computer use in 1997. \medskip
\item Estimate the regression equation:
\[
\def\stacktype{L}
C_{j,1997} = 
\begin{matrix}
-24.56 \\
(19.18)
\end{matrix}
+ 
\begin{matrix}
1.85 \\
(0.48)
\end{matrix}
\times RTS_{j,1960} \tag{12}  \label{eq12}
\]
\item An industry that is 10 percentage points intenser in routine tasks is 19 percentage points higher up in the distribution of computer adoption.
\end{itemize}
\end{frame}

\section{III. Trends in Job Task Input, 1960-1998}

\subsection{III.A Aggregate trends}

\begin{frame}{III. Trends in job task input, 1960-1998}
\begin{enumerate}
\item[\textcolor{red}{III.A}]  \textcolor{red}{Aggregate trends} \bigskip
\item[\textcolor{gray}{III.B}]  \textcolor{gray}{Task changes within and between industries} \bigskip
\end{enumerate}
\end{frame}

\newpage
\begin{center}
\begin{figure}
\includegraphics[height=8.5cm,width=\textwidth]{figure1.png}
\end{figure} 
\end{center}
\newpage

\newpage
\begin{center}
\begin{figure}
\includegraphics[height=8.5cm,width=\textwidth]{figure2.png}
\end{figure} 
\end{center}
\newpage

\subsection{III.B Task changes within and between industries}

\begin{frame}{III. Trends in job task input, 1960-1998}
\begin{enumerate}
\item[\textcolor{gray}{III.A}]  \textcolor{gray}{Aggregate trends} \bigskip
\item[\textcolor{red}{III.B}]  \textcolor{red}{Task changes within and between industries} \bigskip
\end{enumerate}
\end{frame}

\begin{frame}{Task shifts within- and between-industries}
\begin{itemize}
\item The change over time in mean task input $k$ can be decomposed as (dropping time subscripts):
\[
\Delta T_{k}=\sum\nolimits_{j} \Delta E_{j} \bar{\alpha}_{jk}+\sum\nolimits_{j} \Delta \alpha _{jk} \bar{E}_{j} 
\]
with $ \Delta E_{j}$ the change over time in employment in industry $j$; $ \bar{ \alpha}_{jk}$ the mean of task $k$ in industry $j$ averaged across years; $ \Delta \alpha_{jk}$ the change over time in the mean of task $k$ in industry $j$; and $ \bar{E_{j}}$ employment in industry $j$ averaged over time.\bigskip
\item The first term is the between-industry component and the second term is the within-industry component.
\end{itemize}
\end{frame}

\newpage
\begin{center}
\begin{figure}
\includegraphics[height=8.5cm,width=\textwidth]{table2.png}
\end{figure} 
\end{center}
\newpage

\section{IV. Computerization and Task Change: Industry-Level Relationships}

\subsection{IV.A Industry computerization and task trends over four decades}

\begin{frame}{IV. Computerization and task changes within industries}
\begin{enumerate}
\item[\textcolor{red}{IV.A}] \textcolor{red}{Industry computerization and task trends over four decades} \bigskip
\item[\textcolor{gray}{IV.B}] \textcolor{gray}{Using composite DOT variables} \bigskip
\item[\textcolor{gray}{IV.C}] \textcolor{gray}{Employing contemporaneous measures of computer and capital investment}
\end{enumerate}
\end{frame}

\begin{frame}{Industry computerization and task trends, 1960-1998}
\begin{itemize}
\item The model predicts that industries that adopt computers more will reduce routine and increase non-routine labor tasks more.\medskip
\item Estimate:
\[
\Delta T_{jk \tau} = \alpha + \phi \Delta C_{j} + \epsilon_{jk \tau} \tag{13} \label{eq13}
\]
with $ \Delta T_{jk \tau} $ the change in $j$'s input of task $k$ between $t$ and $ \tau $ and $ \Delta C_{j} $ the annual change in the percentage of industry workers using a computer between 1984 and 1997 from CPS. \medskip
\item Estimate this equation for each of the four decades to compare 1960s-1970s (pre-computer era) to 1980s-1990s.
\end{itemize}
\end{frame}

\newpage
\begin{center}
\begin{figure}
\includegraphics[height=8.5cm,width=\textwidth]{table3.png}
\end{figure} 
\end{center}
\newpage

\subsection{IV.B Using composite DOT variables}

\begin{frame}{IV. Computerization and task changes within industries}
\begin{enumerate}
\item[\textcolor{gray}{IV.A}] \textcolor{gray}{Industry computerization and task trends over four decades} \bigskip
\item[\textcolor{red}{IV.B}] \textcolor{red}{Using composite DOT variables} \bigskip
\item[\textcolor{gray}{IV.C}] \textcolor{gray}{Employing contemporaneous measures of computer and capital investment}
\end{enumerate}
\end{frame}

\newpage
\begin{center}
\begin{figure}
\includegraphics[height=8.5cm,width=\textwidth]{appendixtable3.png}
\end{figure} 
\end{center}
\newpage

\subsection{IV.C Employing contemporaneous measures of computer and capital investment}

\begin{frame}{IV. Computerization and task changes within industries}
\begin{enumerate}
\item[\textcolor{gray}{IV.A}] \textcolor{gray}{Industry computerization and task trends over four decades} \bigskip
\item[\textcolor{gray}{IV.B}] \textcolor{gray}{Using composite DOT variables} \bigskip
\item[\textcolor{red}{IV.C}] \textcolor{red}{Employing contemporaneous measures of computer and capital investment}
\end{enumerate}
\end{frame}

\begin{frame}{Using measures of computer and capital investment}
\begin{itemize}
\item The CPS computer use measure only exists for the 1980s and 1990s. Use computer and capital investments for 1959-1998 from the National Income and Product Accounts (NIPA). \medskip
\item Estimate the following stacked first-difference models:
\[
\Delta T_{jk \tau} = \alpha + \delta_{70-80} +\delta_{80-90} + \delta_{90-98} + \varphi CI_{j \tau} + \theta KI_{j \tau} + \epsilon_{jk \tau} \tag{14} \label{eq14}
\]
with $CI$ computer and $KI$ capital investment per FTE. \medskip
\item Cluster standard errors at NIPA sector and decade ($168=42x4$ clusters) to account for aggregation or at NIPA sector ($42$ clusters) to also account for serial correlation.  
\end{itemize}
\end{frame}

\newpage
\begin{center}
\begin{figure}
\includegraphics[height=8.5cm,width=\textwidth]{table4.png}
\end{figure} 
\end{center}
\newpage

\section{V. Task Change Within Education Groups and Occupations}

\subsection{V.A Within-industry task shifts by education group: 1980-1998}

\begin{frame}{V. Task change within education groups and occupations}
\begin{enumerate}
\item[\textcolor{red}{V.A}] \textcolor{red}{Within-industry task shifts by education group: 1980-1998}  \bigskip
\item[\textcolor{gray}{V.B}]  \textcolor{gray}{Task shifts within occupations}
\end{enumerate}
\end{frame}

\begin{frame}{Task changes within education groups}
\begin{itemize}
\item Assume that educated workers have a comparative advantage in doing non-routine labor tasks:
\[
\eta_{i} \equiv n_{i}/r_{i} \in (0, \infty) \text{ is higher for educated workers}
\]
\item Shifts towards non-routine tasks against routine labor tasks could then be explained by an increase in the relative demand for educated workers due to computerization.\medskip
\item To show that changes in labor task demands are an underlying cause and not a symptom of changes in relative demand for skills, we look at labor task changes within education groups.
\end{itemize}
\end{frame}

\begin{frame}{Within-industry task shifts by education}
\begin{itemize}
\item We can test whether industries change task assignments of workers with given levels of education.\medskip
\item Estimate equation (13) by education group:
\[
\Delta T_{ijk \tau} = \alpha_{i} + \phi_{i} \Delta C_{j} + \epsilon_{ijk \tau} \tag{15} \label{eq15}
\]
with $i$ one of four education groups. \medskip
\item Use industry task measures for 1980-1998 and (almost) contemporaneous computer use data for 1984-1997.
\end{itemize}
\end{frame}

\newpage
\begin{center}
\begin{figure}
\includegraphics[height=8.5cm,width=\textwidth]{table5a.png}
\end{figure} 
\end{center}
\newpage

\newpage
\begin{center}
\begin{figure}
\includegraphics[height=8.5cm,width=\textwidth]{table5b.png}
\end{figure} 
\end{center}
\newpage

\begin{frame}{Within-industry task shifts by education (Table V)}
\begin{itemize}
\item Industry-level computerization is strongly predictive of shifts towards non-routine and against routine labor tasks within education groups.\medskip
\item A decomposition of industry task changes into within and between education groups shows that the within component is quantitatively important.\medskip
\item This suggests that labor task changes due to computerization are an antecedent to educational upgrading, rather than a symptom of it.
\end{itemize}
\end{frame}

\subsection{V.B Task shifts within occupations}

\begin{frame}{V. Task change within education groups and occupations}
\begin{enumerate}
\item[\textcolor{gray}{V.A}] \textcolor{gray}{Within-industry task shifts by education group: 1980-1998}  \bigskip
\item[\textcolor{red}{V.B}]  \textcolor{red}{Task shifts within occupations}
\end{enumerate}
\end{frame}

\begin{frame}{Task shifts within occupations}
\begin{itemize}
\item Task measures were taken from 1977 DOT and are occupation specific.\medskip
\item This assumes that tasks done within occupations are static, which is an unrealistic assumption in the long-run.\medskip
\item In particular, computerization is expected to reduce routine and increase non-routine labor tasks also within occupations. \medskip
\item We can analyse within-occupation task changes by comparing the 1977 and 1991 DOT measures. 
\end{itemize}
\end{frame}

\begin{frame}{Description of the job of ``Secretary''}
\begin{block}{Occupation Outlook Handbook, 1976}
\textit{``Secretaries relieve their employers of routine duties so they can work on more important matters. Although most secretaries type, take shorthand, and deal with callers, the time spent on these duties varies in different types of organizations."}
\end{block}
\begin{block}{Occupation Outlook Handbook, 2000}
\textit{``As technology continues to expand in offices, the role of the secretary has greatly evolved. Office automation and organizational restructuring have led secretaries to assume a wide range of new responsibilities once reserved for managerial and professional staff. Many secretaries now provide training and orientation to new staff, conduct research on the internet, and learn to operate new office technologies."}
\end{block}
\end{frame}

\begin{frame}{Task shifts within occupations}
\begin{itemize}
\item Merge 1977 DOT with 1991 DOT to capture changes in task content within occupations. \medskip
\item Estimate the following equation:
\[
\Delta T_{mk} = \alpha + \xi \Delta C_{m} + \epsilon_{mk} \tag{16} \label{eq16}
\]
with $\Delta T_{mk} $ the change in occupation $m$ of task $k$ between 1977 and 1991. \medskip
\item Note that variation in $\Delta T_{mk} $ stems exclusively from DOT examiners' reevaluation of the task content of individual occupations between 1977 and 1991.
\end{itemize}
\end{frame}

\newpage
\begin{center}
\begin{figure}
\includegraphics[height=8.5cm,width=\textwidth]{table6.png}
\end{figure} 
\end{center}
\newpage

\begin{frame}{Task shifts within occupations (Table VI)}
\begin{itemize}
\item Occupations with more computerization had stronger shifts towards non-routine and against routine labor tasks.\bigskip
\item This is not driven by shifts in educational and gender distributions of workers within occupations. \bigskip
\item In net, shifts in job contents are a pervasive feature of the data and are concentrated in industries and occupations that adopted computer technology more rapidly.
\end{itemize}
\end{frame}

\section{VI. Quantifying the Magnitude of Task Structure Changes}

\begin{frame}{Quantifying the magnitude of task changes}
\begin{itemize}
\item The economic significance of task shifts is not clear given that they do not have a familiar scale. \medskip
\item To meaningfully quantify task changes, calculate their contribution to changes in the demand for more-educated relative to less-educated workers between 1970 and 1998. \medskip
\item To do this, three steps are taken.
\end{itemize}
\end{frame}

\begin{frame}{Quantifying the magnitude of task changes}
\begin{enumerate}
\item Estimate a fixed-coefficient model of educational requirements in industries and occupations as a function of their task inputs:
\[
\text{College share}_{j} = \alpha + \sum_{k=1}^{4} \pi_{k} T_{j}^{k} + \epsilon_{j} \tag{17} \label{eq17}
\]
\item Use $ \hat{\pi}_{k}$ to predict changes in the aggregate college share:
\[
\hat{\Delta} \text{College share}_{70-98} = \alpha + \sum_{k=1}^{4} \hat{\pi}_{k} \Delta T_{70-98}^{k} \tag{18} \label{eq18}
\]
with $\Delta T_{70-98}^{k}$ observed or predicted from (13), (14) or (16). \medskip
\item Compare predicted changes with predicted demand changes from the canonical model.
\end{enumerate}
\end{frame}

\newpage
\begin{center}
\begin{figure}
\includegraphics[height=8.5cm,width=\textwidth]{table7a.png}
\end{figure} 
\end{center}
\newpage

\newpage
\begin{center}
\begin{figure}
\includegraphics[height=8.5cm,width=\textwidth]{table7b.png}
\end{figure} 
\end{center}
\newpage

\begin{frame}{Quantifying the magnitude of task changes (Table VII)}
\begin{itemize}
\item For 1980-1998, task changes due to computerization explain 60\% to 90\% of the estimated increase in the relative demand for college employment. \medskip
\item 60\% of the computer contribution to rising educational demand is due to shifts in the task composition between occupations. \medskip
\item In net, this illustrates that changes in task demands due to computerization are economically large.  
\end{itemize}
\end{frame}

\section{VII. Conclusions}

\begin{frame}{Conclusions}
\begin{itemize}
\item This paper develops a simple model of how computers substitute for workers performing routine tasks and complements workers doing non-routine tasks. \medskip
\item The model also predicts that industries intensive in routine labor tasks make larger investments in computers. \medskip
\item Using DOT data, these predictions are confirmed across several margins of task changes. \medskip
\item These task shifts are an important contributing factor underlying the increase in the relative demand for educated workers in the 1980s and 1990s.
\end{itemize}
\end{frame}

\end{document}
